\section{Evaluation}
\label{sec:experiments}

The evaluation separates properties that can be established without a model
provider from observations that require real semantic-verifier calls. We ask
three questions. \textbf{RQ1} tests whether the implemented evidence contract
accepts contract-valid reports and intercepts controlled integrity, grounding,
scope, and conflict defects. \textbf{RQ2} tests whether SpanVouch preserves
determinism, recovery, dependency boundaries, and delivery behavior across its
full release gate. \textbf{RQ3} tests whether the multi-domain,
multi-framework B0--B5 artifact pipeline executes end to end without leaking
labels or requiring a live provider. The v0.7 formal evidence is a separate
DeepSeek-only observation of the declared matrix. Its claim gates remain
fail-closed and unresolved; no outcome is promoted to a hypothesis support
claim.

\subsection{Frozen Deterministic Benchmark}

The primary controlled benchmark contains 36 SupportLab diagnosis candidates
derived from 20 frozen traces. Twenty candidates are contract-valid,
unmodified reports. The other 16 contain one programmed defect from six
families: unsupported scope (6), invalid selector (2), evidence-value hash
mismatch (2), ungrounded claim (2), ungrounded critical span (2), and diagnosis
conflict (2). The deterministic verifier processes every candidate without a
provider call. A valid candidate succeeds only when it receives
\textsc{verified}; a mutation is intercepted when it receives
\textsc{needs-evidence} or \textsc{review-required} with the corresponding
typed finding.

The checked-in reference bundle binds the candidate and label manifests,
configuration, code revision, runtime, and output hashes. Its artifact-manifest
SHA-256 is
\nolinkurl{f8f33da259ff055edf626ac50f35102e1763e7ee00e8d9a046003e945f1ee3a5};
the metrics SHA-256 is
\nolinkurl{c3fc1f4fc2015cbc0a3d6691bf01da98e1a77f7ac139d37f7e75cd2038742f9b}.

\subsection{Release and Recovery Validation}

The release gate combines unit, contract, architecture, integration, process,
and end-to-end tests. Contract tests cover the six public schema roots and
valid/invalid fixtures. Architecture tests enforce dependency direction and
provider-label isolation. Recovery tests exercise idempotency, compare-and-swap
updates, owner-fenced leases, multi-process SQLite access, and state persistence
across an API-container restart. Static and delivery gates run Ruff, strict
mypy, locked wheel construction, Compose validation, and an unprivileged
container smoke test.

\subsection{Offline Multi-Framework Matrix}

The pipeline validation crosses two domains (SupportLab and OpsLab), two
frameworks (LangGraph and AutoGen), and six verifier conditions. B0 accepts a
contract-valid diagnosis without a verifier; B1 applies deterministic
verification; B2 and B3 exercise shared-context and isolated semantic-verifier
interfaces; B4 exercises an isolated cross-model interface; and B5 composes
deterministic and isolated semantic interfaces. In this validation, semantic
conditions use declared fake providers and therefore carry no semantic-
effectiveness interpretation.

The run schedules $2\times2\times6=24$ cells, records four framework-adapter
executions and two parity comparisons, joins labels only after the provider
phase, and generates metrics, claim gates, paper assets, and a manifest. The
artifact explicitly records \code{fake\_\allowbreak{}evidence=true}, zero provider calls,
and zero GPU calls. Repeating the run must produce byte-identical bundles.

\subsection{Formal DeepSeek-only Evidence and Claim--Evidence Map}

The public aggregate bundle in
\nolinkurl{evals/reports/reference/phase5-formal-deepseek-only/} binds the
analysis input, evaluated-results manifest, and every copied result asset. The
evaluated-results SHA-256 is
\nolinkurl{bc09f1b134de9370a3b5209fa5e959bce01abbcdf05c8456af1f069fc4cd3088}.
The matrix scheduled and evaluated 2,148 plans with zero missing cells. B0--B3
completed under DeepSeek-only conditions; B4 and B5 were \textsc{policy-skipped}
and have no Qwen results.

\begin{table}[t]
    \caption{v0.7 claim--evidence map. Entries are observational and remain
    bounded by the stated artifact and scope.}
    \label{tab:claim-evidence-map}
    \centering
    \scriptsize
    \begin{tabular}{@{}p{0.26\columnwidth}p{0.36\columnwidth}p{0.30\columnwidth}@{}}
        \toprule
        Claim surface & Evidence & Permitted interpretation \\
        \midrule
        Matrix completion & evaluated-results manifest & 2,148 evaluated; 0 missing \\
        SupportLab B3 vs B2 & paired-effects.csv & observed risk--coverage tradeoff \\
        OpsLab B3 vs B2 & paired-effects.csv & preliminary observed replication \\
        H1--H5 & claim-gates.json & unresolved; no support claim \\
        B4/B5 & metrics and claim assets & policy-skipped; no Qwen result \\
        \bottomrule
    \end{tabular}
\end{table}

\subsection{Real-Provider and Risk Protocol}

The result-bearing semantic study uses the same frozen traces and diagnoses
across B0--B5. Provider-visible inputs exclude labels, mutation metadata,
expected findings, split identity, credentials, and hidden reasoning. The
primary semantic endpoint is a correct, evidence-supported diagnosis under
execution-derived labels; paired bootstrap intervals resample scenario-template
clusters and prespecified primary comparisons use Holm correction.

The group-selective study freezes development, calibration, and untouched test
groups before outcome analysis. It reports every $(K_\lambda^G,M_\lambda^G)$
pair, simultaneous upper bound, accepted-group minimum, feasibility decision,
and selected candidate. Provider failures and contract-invalid outputs remain
visible. No value from the fake-provider matrix enters this study's
effectiveness or risk calculations. The v0.7 DeepSeek-only observations do not
establish the preregistered H1--H5 hypotheses, and no empirical target-risk
certificate is reported.
